\scanpage{7}
\mspage{(5)}
car dans $S\{X_i\}[dX_i]$ l'idéal formé des formes \inserted{à p.d.} d'augmentation $\in J$ est à p.d.
\[
\left(S\{X_i\}_{0\leq i\leq n}/(\textstyle\sum X_i-t)_{\mathrm{pd}}\right)\otimes_S\bigwedge^*\left(S^{[n]}/\operatorname{diag}\unsure{S^{[n]}}\right)
\]
C'est une $S$-algèbre diff. \inserted{graduée} anti comm., à degrés $\geq0$, \inserted{augmentée vers $S/J$} à p.d. sur l'idéal \deleted{engendré par $\bigwedge^+$, $S\{X_i\}^+$ et $J\{X_i\}$ \unsure{dans le premier}} l'idéal noyau de l'augmentation
\[
C^\bullet_{\mathrm{DRpd}\,[n]}(S,J,t)\longrightarrow S/J\ ;
\]
\margin{la diff est compatible avec les p.d. i.e. $d(x^{(n)})=x^{(n-1)}dx$}
et comme telle isomorphe à $S\{X_0,\ldots,X_{n-1}\}[dX_0,\ldots,dX_{n-1}]$ qui \uline{est une résolution de $S$}.
\deleted{Donc a) est acquis. Mais faute de b), on ne peut dire a priori ce \unsure{(et il ne sera)} que fait $H^{\unsure{q}}C^\bullet$.}
Pour $[n]$ variable, on trouve
\[
C^\bullet_{\mathrm{DRpd}\,*}{}^{(S,J,t)}=\left(C^\bullet_{\mathrm{DRpd}\,[n]}{}^{(S,J,t)}\right)_{n\geq0}
\]
qui est une rés. $\frac12$ simpliciale (et $=$ simpliciale) de $S$, \inserted{avec aug. vers $S/J=k$, et p.d. sur l'idéal d'augm., comp. à d}. Elle dépend fonctoriellement de $(S,J,t)$, et elle permet de former, \inserted{pour $\mathcal X_*\in(\mathrm{Ss})$}
\[
C^\bullet_{\mathrm{DRpd}}(\mathcal X_*;S,J,t)=\operatorname{Hom}(\mathcal X_*,C^\bullet_{\mathrm{DRpd}\,*}(S,J,t))
\]
foncteur contravariant en $\mathcal X_*$ ! (et si $X$ esp. top.
\[
\begin{aligned}
C^\bullet_{\mathrm{DRpd}}(X;\quad)&=C^\bullet_{\mathrm{DRpd}}(S_*(X);\text{---})\\
&=\operatorname{Hom}(\unsure{C_*(X)|\mathbb Z_*(X)},C^\bullet_{\mathrm{DRpd}\,*}(\quad))
\end{aligned}
\]
\margin{à valeurs dans les $\unsure{S}$-algèbres graduées diff. $S/J$-augmentées, à p.d. sur l'idéal d'augmentation, compatibles avec la diff.}
Mais on ne peut dire en général quelle est sa cohomologie, \inserted{(on a seulement $H^{\unsure{\bullet|0}}_{\mathrm{DRpd}}(\mathcal X_*;S,J,t)\longrightarrow H^{\unsure{\bullet|0}}(\mathcal X_*,S)$)} et en tous cas elle ne sera pas isom. (pour $\mathcal X_*$ variable) à $H^\bullet(\mathcal X_*,S)$.
\mspage{6}
